
\needspace{4\baselineskip}
\color{uniblue}{\subsection{Защита от обрыва провода (ЗОП)}}
\color{black}

\begin{enumerate}[label=\arabic{section}.\arabic{subsection}.\arabic*, labelsep=4pt, leftmargin=0pt, itemindent=57pt]

\item
Защита от обрыва провода предназначена для обнаружения обрыва токоведущего проводника или выявления несимметричных режимов работы и неисправности измерительных трансформаторов тока. Принцип действия защиты основан на оценке соотношения токов обратной и прямой последовательностей или на оценке значения тока обратной последовательности.

\item
В устройстве защита от обрыва провода представлена функцией «ЗОП».

\item
Основные характеристики функции:
\begin{itemize}
    \item погрешность измерительного органа несимметрии не более 3\%;
    \item время срабатывания измерительного органа прямой последовательности при значении уставки $K_{\textsl{несим}}=1$, значении тока обратной последовательности равном $0,2\cdot I_{\textsl{ном}}$ и повышении тока прямой последовательности скачком от 0 до $0,08\cdot I_{\textsl{ном}}$ не более 40 мс;
    \item время возврата измерительного органа прямой последовательности при значении уставки $K_{\textsl{несим}}=1$, значении тока обратной последовательности равном $0,2\cdot I_{\textsl{ном}}$ и снижении тока прямой последовательности скачком от $0,08\cdot I_{\textsl{ном}}$ до 0 не более 40 мс;
    \item время срабатывания измерительного органа по несимметрии при значении уставки $K_{\textsl{несим}}=1$, значении тока прямой последовательности равном $0,08\cdot I_{\textsl{ном}}$ и повышении тока обратной последовательности скачком от 0 до $0,16\cdot I_{\textsl{ном}}$ не более 40 мс;
    \item время возврата измерительного органа по несимметрии при значении уставки $K_{\textsl{несим}}=1$, значении тока прямой последовательности равном $0,08\cdot I_{\textsl{ном}}$ и снижении тока обратной последовательности скачком от $0,16\cdot I_{\textsl{ном}}$ до 0 не более 40 мс.
\end{itemize}

\item
В качестве входных значений функция использует значения токов прямой и обратной последовательностей ($I_1$ и $I_2$). Расчет $I_1$ и $I_2$ выполняется по следующим формулам:

% ФОРМУЛА №1
\begin{equation*}
    %\label{eq:zop1}
    \dot{I}_1=\frac{1}{3}\left[\dot{I}_A+\dot{I}_Be^{j120}+\dot{I}_Ce^{j240}\right],
\end{equation*}
\begin{equation*}
    %\label{eq:zop2}
    \dot{I}_2=\frac{1}{3}\left[\dot{I}_A+\dot{I}_Be^{j240}+\dot{I}_Ce^{j120}\right],
\end{equation*}
\begin{eqexpl}[25mm]
%\item{$\dot{I}_1$} вычисленный ток прямой последовательности;
%\item{$\dot{I}_2$} вычисленный ток обратной последовательности;
\item{$\dot{I}_A, \dot{I}_B, \dot{I}_C$} фазные токи стороны ВН трансформатора.
\end{eqexpl}

В зависимости от выбранного режима работы функция сравнивает вычисленные значения тока $I_2$ или значение отношения $I_2/I_1$ с заданной уставкой.

\item
Ввод функции в работу на этапе параметрирования устройства осуществляется своим программным переключателем SGF1 «Ввод\_функции». Информация о введенной функции отображается выходным сигналом «ЗОП: Ввод».
\item
Функция может быть выведена из работы оперативно путем активации сигнала «ОВ~ЗОП», что характеризуется активным состоянием сигнала «ЗОП: Оперативный вывод» на выходе функции.
\item
Функция может контролировать, в зависимости от выбранного состояния программного переключателя SGF2 «Реж\_ЗОП»:
\begin{itemize}
 \item повышение тока обратной последовательности (\textsl{ИО I2>}); 
 \item повышение величины отношения тока обратной последовательности к току прямой последовательности (\textsl{ИО I2/I1>});
 \item повышение тока обратной последовательности (\textsl{ИО I2>}) и повышение величины отношения тока обратной последовательности к току прямой последовательности (\textsl{ИО I2/I1>}).
\end{itemize}
 При превышении заданной уставки по перечисленным выше ИО активируется сигнал «ЗОП: Пуск», далее через выдержку времени таймера T1 (уставка «Тср») выдается сигнал срабатывания «ЗОП: Срабатывание сигн». Сигнал отключения от ЗОП «ЗОП: Срабатывание» формируется одновременно с сигналом «ЗОП: Срабатывание сигн» при неактивном состоянии входного сигнала «ЗОП на сигн». ИО I2/I1 вводится автоматически в работу при превышении тока прямой последовательности величины 0,04Iном.
\item
Алгоритм работы функции выполнен в соответствии с рисунком \ref{zop:img1}. В таблице \ref{zop:tbl1} приведены параметры, необходимые для настройки функции.

\vspace{3mm}
\begin{figure}[H]
\centering
\includegraphics[width=1\textwidth,height=1\textheight,keepaspectratio]{\fbpath/55. ЗОП/_latex/img/zop.pdf}
\captionof{figure}{Функциональная схема логики ЗОП}\label{zop:img1}
\end{figure}


\end{enumerate}
\FloatBarrier % Все float'ы ДО этой строки будут размещены здесь.